\documentclass[12pt,a4paper]{amsart}
\usepackage{amsmath,amssymb,amsthm}
\usepackage{mathtools}
\usepackage[utf8]{inputenc}
\usepackage[T1]{fontenc}
\usepackage{hyperref}
\usepackage{enumitem,booktabs,array}
\usepackage{geometry}
\geometry{margin=2.5cm}

\newtheorem{theorem}{Theorem}[section]
\newtheorem{lemma}[theorem]{Lemma}
\newtheorem{corollary}[theorem]{Corollary}
\newtheorem{proposition}[theorem]{Proposition}
\newtheorem{conjecture}[theorem]{Conjecture}
\theoremstyle{definition}
\newtheorem{definition}[theorem]{Definition}
\newtheorem{example}[theorem]{Example}
\newtheorem{remark}[theorem]{Remark}

\author{Michael Fuhrmann}
\address{Specialist Mathematics Project, Build 143}

\begin{document}
\title{Theta Functions: Bridges Between Analysis, Number Theory, and Geometry}

\begin{abstract}
Theta functions stand at the crossroads of classical analysis, number theory, and
algebraic geometry. In this paper we develop the theory of Jacobi theta functions
$\vartheta_1, \vartheta_2, \vartheta_3, \vartheta_4$, establish their transformation
formulas under the action of the modular group $\mathrm{SL}_2(\mathbb{Z})$, and
interpret them as modular forms of half-integral weight. We give complete proofs
of Jacobi's four-square theorem
$r_4(n) = 8\sum_{d \mid n,\, 4 \nmid d} d$,
the Poisson summation approach to the functional equation of the Riemann zeta function,
and the classical Rogers--Ramanujan identities. The Dedekind eta function $\eta(\tau)$
is introduced and linked to the discriminant $\Delta = \eta^{24}$.
Higher-dimensional Siegel theta functions and the uniformisation of elliptic curves
via Weierstra\ss{} $\wp$-functions are surveyed. Finally, we discuss Zwegers'
celebrated 2002 result that places Ramanujan's mock theta functions within the
framework of harmonic Maass forms, thereby resolving a mystery that had persisted
for more than eighty years.
\end{abstract}

\maketitle

\tableofcontents

%% ============================================================
\section{Introduction}
%% ============================================================

Since their introduction by Carl Gustav Jacob Jacobi in his
\textit{Fundamenta Nova} of 1829 \cite{Jacobi1829},
theta functions have occupied a central position in mathematics.
At first sight they appear as elegant $q$-series;
their deeper significance, however, lies in the surprising number of
fundamental theorems they encode.

A single formula,
\begin{equation}\label{eq:theta3basic}
  \vartheta_3(\tau) \;=\; \sum_{n=-\infty}^{\infty} q^{n^2},
  \qquad q = e^{\pi i \tau},\quad \operatorname{Im}(\tau)>0,
\end{equation}
connects:
\begin{itemize}[leftmargin=2em]
  \item \textbf{Analytic number theory}: the coefficient of $q^n$ in
        $\vartheta_3(\tau)^k$ counts the number of representations
        $r_k(n)$ of $n$ as a sum of $k$ integer squares;
  \item \textbf{Complex analysis}: the transformation
        $\vartheta_3(\tau) \mapsto \sqrt{-i\tau}\,\vartheta_3(\tau)$
        under $\tau \mapsto -1/\tau$ is equivalent to the Poisson
        summation formula and yields the functional equation of
        Riemann's zeta function;
  \item \textbf{Modular forms}: $\vartheta_3(\tau)^2$ is a modular form
        of weight $1$ for the congruence subgroup $\Gamma_0(4)$;
  \item \textbf{Algebraic geometry}: every elliptic curve over $\mathbb{C}$
        is uniformised by $\mathbb{C}/\Lambda$ and the corresponding
        lattice theta series encodes arithmetic data about the curve.
\end{itemize}

The paper is organised as follows.
Section~\ref{sec:jacobi} introduces all four Jacobi theta functions.
Section~\ref{sec:transformation} establishes the transformation formulas.
Section~\ref{sec:eta} treats the Dedekind eta function.
Section~\ref{sec:modular} interprets theta functions as modular forms.
Section~\ref{sec:quadratic} applies them to representations by quadratic forms.
Section~\ref{sec:poisson} derives the functional equation of $\zeta(s)$.
Section~\ref{sec:siegel} surveys Siegel theta functions.
Section~\ref{sec:elliptic} connects theta functions to elliptic curves.
Section~\ref{sec:ramanujan} discusses Ramanujan's identities and mock theta functions.

\medskip
\noindent\textbf{Notation.}
Throughout, $\tau \in \mathfrak{H} = \{\tau \in \mathbb{C} : \operatorname{Im}(\tau)>0\}$
and $q = e^{2\pi i\tau}$ unless stated otherwise (for Jacobi's own convention
$q_J = e^{\pi i \tau}$ we write $q_J$).
We write $e(x) = e^{2\pi i x}$.

%% ============================================================
\section{Jacobi Theta Functions}\label{sec:jacobi}
%% ============================================================

\subsection{Definitions as $q$-series}

\begin{definition}[Jacobi Theta Functions]\label{def:jacobi}
For $z \in \mathbb{C}$ and $\tau \in \mathfrak{H}$, set $q = e^{\pi i \tau}$
(Jacobi's convention). The four Jacobi theta functions are:
\begin{align}
  \vartheta_1(z|\tau) &= -i \sum_{n=-\infty}^{\infty}
    (-1)^n\, q^{(n+\tfrac{1}{2})^2}\, e^{(2n+1)\pi iz},
    \label{eq:theta1}\\
  \vartheta_2(z|\tau) &= \sum_{n=-\infty}^{\infty}
    q^{(n+\tfrac{1}{2})^2}\, e^{(2n+1)\pi iz},
    \label{eq:theta2}\\
  \vartheta_3(z|\tau) &= \sum_{n=-\infty}^{\infty}
    q^{n^2}\, e^{2n\pi iz},
    \label{eq:theta3}\\
  \vartheta_4(z|\tau) &= \sum_{n=-\infty}^{\infty}
    (-1)^n\, q^{n^2}\, e^{2n\pi iz}.
    \label{eq:theta4}
\end{align}
\end{definition}

\begin{remark}
Using $q = e^{\pi i\tau}$ (so $|q| = e^{-\pi\operatorname{Im}\tau} < 1$),
the half-integer exponents $(n+\frac12)^2$ appear naturally.
Many modern texts prefer $q = e^{2\pi i\tau}$; in that convention
$\vartheta_3(\tau) = \sum_{n\in\mathbb{Z}} q^{n^2/2}$,
which shifts the weight by a factor of $\sqrt{2}$ in the transformation formula.
We follow Jacobi's original convention in this section and the standard
$q = e^{2\pi i\tau}$ everywhere else.
\end{remark}

\subsection{Convergence}

\begin{proposition}[Uniform convergence]\label{prop:convergence}
Each of the four series in Definition~\ref{def:jacobi} converges absolutely and
uniformly on every compact subset of $\mathbb{C} \times \mathfrak{H}$.
\end{proposition}

\begin{proof}
We estimate the general term of $\vartheta_3$.  For fixed $\operatorname{Im}(\tau) \geq \delta > 0$
and $|z| \leq R$ one has
\[
  \bigl|q^{n^2} e^{2n\pi iz}\bigr|
  = e^{-\pi n^2 \operatorname{Im}(\tau)} \cdot e^{-2\pi n \operatorname{Im}(z)}
  \leq e^{-\pi \delta n^2 + 2\pi |n| R}.
\]
For $|n|$ large the term $\pi\delta n^2$ dominates $2\pi|n|R$, so the series is bounded
by a convergent series independent of $(z,\tau)$ in the compact set.
The argument for $\vartheta_1, \vartheta_2, \vartheta_4$ is identical.
\end{proof}

\subsection{Product representations}

The Jacobi triple product identity (Theorem~\ref{thm:triple}) gives product
representations for all four theta functions:
\begin{align}
  \vartheta_3(z|\tau) &= \prod_{n=1}^{\infty}
    (1-q^{2n})(1+q^{2n-1}e^{2\pi iz})(1+q^{2n-1}e^{-2\pi iz}).
\end{align}

\begin{theorem}[Jacobi Triple Product, 1829]\label{thm:triple}
For $|q|<1$ and $z \neq 0$,
\[
  \sum_{n=-\infty}^{\infty} z^n q^{n^2}
  = \prod_{m=1}^{\infty}(1-q^{2m})(1+zq^{2m-1})(1+z^{-1}q^{2m-1}).
\]
\end{theorem}

\begin{proof}[Proof sketch]
Let $f(z) = \prod_{m\geq1}(1+zq^{2m-1})(1+z^{-1}q^{2m-1})$.
One shows $f(zq^2) = z^{-1}q^{-1}f(z)$, so $f$ has the same quasi-periodicity
as $\sum_n z^n q^{n^2}$.  Comparing coefficients (Laurent expansion in $z$)
and using the infinite product $\prod(1-q^{2m})$ for normalisation yields the identity.
A complete proof by induction on the number of factors appears in \cite{WhittakerWatson1927}.
\end{proof}

%% ============================================================
\section{Transformation Formulas}\label{sec:transformation}
%% ============================================================

\subsection{Periodicity}

\begin{proposition}[Periodicity in $z$]\label{prop:period_z}
The theta functions satisfy:
\begin{align}
  \vartheta_3(z+1|\tau) &= \vartheta_3(z|\tau),\\
  \vartheta_3(z+\tau|\tau) &= q^{-1} e^{-2\pi iz}\, \vartheta_3(z|\tau),
\end{align}
and similarly for $\vartheta_1, \vartheta_2, \vartheta_4$.
\end{proposition}

\begin{proof}
Replacing $n \mapsto n$ when $z \mapsto z+1$ leaves the exponent $2n\pi iz$ unchanged
modulo $2\pi in$, which yields $e^{2\pi in} = 1$.
For $z \mapsto z + \tau$: the factor $e^{2n\pi i(z+\tau)} = e^{2n\pi iz}\cdot q^{2n}$,
while the $q^{n^2}$ factor becomes $q^{n^2}$.  Hence
$q^{n^2} e^{2n\pi i(z+\tau)} = q^{(n+1)^2 - 1} e^{2\pi iz \cdot 2n}
= q^{-1}e^{-2\pi iz} \cdot q^{(n+1)^2} e^{2(n+1)\pi iz}$,
giving the stated quasi-periodicity.
\end{proof}

\subsection{Modular transformations}

\begin{theorem}[Transformation under $S: \tau \mapsto -1/\tau$]\label{thm:S_transform}
\[
  \vartheta_3\!\left(z \,\Big|\, -\tfrac{1}{\tau}\right)
  = \sqrt{-i\tau}\; e^{\pi i \tau z^2}\; \vartheta_3(z\tau \mid \tau).
\]
In particular, at $z = 0$:
\begin{equation}\label{eq:theta3_S}
  \vartheta_3\!\left(0 \,\Big|\, -\tfrac{1}{\tau}\right)
  = \sqrt{-i\tau}\; \vartheta_3(0|\tau).
\end{equation}
\end{theorem}

\begin{proof}
We apply the Poisson summation formula (Theorem~\ref{thm:poisson}) to the function
$f(x) = e^{\pi i\tau x^2 + 2\pi i z x}$.  Its Fourier transform is
$\hat{f}(\xi) = \tfrac{1}{\sqrt{-i\tau}} e^{\pi i(\xi-z)^2/\tau}$.
Summing over $x = n \in \mathbb{Z}$ on the left and $\xi = n$ on the right yields
exactly \eqref{eq:theta3_S} (after setting $z=0$).
For the general $z$ version an additional shift $x \mapsto x + z$ is applied.
\end{proof}

\begin{theorem}[Transformation under $T: \tau \mapsto \tau+1$]\label{thm:T_transform}
\[
  \vartheta_3(z|\tau+1) = \vartheta_4(z|\tau),
  \qquad
  \vartheta_2(z|\tau+1) = e^{\pi i/4}\, \vartheta_2(z|\tau),
  \qquad
  \vartheta_4(z|\tau+1) = \vartheta_3(z|\tau).
\]
\end{theorem}

\begin{proof}
In the series for $\vartheta_3$ replace $\tau \to \tau+1$:
$q^{n^2} \mapsto e^{\pi i n^2(\tau+1)} = e^{\pi i n^2 \tau} \cdot e^{\pi i n^2}$.
Since $e^{\pi i n^2} = (-1)^n$, we obtain $\vartheta_4$.
The other relations follow analogously.
\end{proof}

\begin{corollary}[Even symmetry]\label{cor:even}
$\vartheta_3(-z|\tau) = \vartheta_3(z|\tau)$, \quad
$\vartheta_4(-z|\tau) = \vartheta_4(z|\tau)$.
\end{corollary}

%% ============================================================
\section{The Dedekind Eta Function}\label{sec:eta}
%% ============================================================

\begin{definition}[Dedekind Eta]\label{def:eta}
For $\tau \in \mathfrak{H}$ and $q = e^{2\pi i\tau}$,
\[
  \eta(\tau) \;=\; q^{1/24}\prod_{n=1}^{\infty}(1-q^n).
\]
\end{definition}

\begin{theorem}[Transformation formulas for $\eta$]\label{thm:eta_transform}
\begin{align}
  \eta(\tau+1) &= e^{\pi i/12}\,\eta(\tau),\label{eq:eta_T}\\
  \eta\!\left(-\tfrac{1}{\tau}\right) &= \sqrt{-i\tau}\;\eta(\tau).\label{eq:eta_S}
\end{align}
\end{theorem}

\begin{proof}[Proof of \eqref{eq:eta_T}]
Replace $\tau \to \tau+1$: $q \to e^{2\pi i(\tau+1)} = q\cdot e^{2\pi i} = q$.
The prefactor $q^{1/24}$ gains a factor $e^{2\pi i/24} = e^{\pi i/12}$.
The product $\prod(1-q^n)$ is unchanged.
\end{proof}

\begin{proof}[Proof of \eqref{eq:eta_S} -- sketch]
One uses the product formula together with Euler's pentagonal number theorem
$\prod_{n\geq1}(1-q^n) = \sum_{n\in\mathbb{Z}}(-1)^n q^{n(3n-1)/2}$,
then applies Poisson summation to the associated theta-type series.
The identity \eqref{eq:eta_S} was first proved by Dedekind (1877) using
Dirichlet series methods; a modern proof via the Poisson formula is in \cite{Apostol1990}.
\end{proof}

\begin{theorem}[Relation to discriminant]\label{thm:delta}
The Ramanujan delta function
\[
  \Delta(\tau) = q\prod_{n=1}^{\infty}(1-q^n)^{24} = \sum_{n=1}^{\infty}\tau(n)\,q^n,
  \qquad q = e^{2\pi i\tau},
\]
satisfies $\Delta(\tau) = \eta(\tau)^{24}$.
In particular, $\Delta$ is a cusp form of weight $12$ for $\mathrm{SL}_2(\mathbb{Z})$.
\end{theorem}

\begin{proof}
Direct computation: $\eta(\tau)^{24} = q^{24/24}\prod(1-q^n)^{24} = q\prod(1-q^n)^{24} = \Delta(\tau)$.
The modular transformation $\Delta((-1/\tau)) = \tau^{12}\Delta(\tau)$ follows from
$\eta(-1/\tau)^{24} = (\sqrt{-i\tau})^{24}\eta(\tau)^{24} = (-i\tau)^{12}\eta(\tau)^{24}$
and $(-i)^{12} = 1$.
\end{proof}

\begin{remark}
Euler's pentagonal number theorem, a special case of Jacobi's triple product,
gives the explicit expansion:
$\prod_{n=1}^{\infty}(1-q^n) = \sum_{k=-\infty}^{\infty}(-1)^k q^{k(3k-1)/2}$,
with exponents $0,1,2,5,7,12,15,\ldots$ (the generalised pentagonal numbers).
\end{remark}

%% ============================================================
\section{Theta Functions as Modular Forms}\label{sec:modular}
%% ============================================================

\subsection{Weight of theta functions}

\begin{definition}[Modular form of half-integral weight]
A holomorphic function $f : \mathfrak{H} \to \mathbb{C}$ is a \emph{modular form
of weight $k/2$} ($k$ odd, positive) for $\Gamma_0(4)$ with character $\chi$ if
\[
  f\!\left(\frac{a\tau+b}{c\tau+d}\right)
  = \chi(d)\,\epsilon_d^{-k}\,(c\tau+d)^{k/2}\,f(\tau)
\]
for all $\begin{pmatrix}a&b\\c&d\end{pmatrix}\in\Gamma_0(4)$, where
$\epsilon_d = 1$ if $d\equiv 1\pmod{4}$ and $\epsilon_d = i$ if $d\equiv 3\pmod{4}$.
\end{definition}

\begin{theorem}[$\vartheta_3^2$ is a modular form of weight $1$]\label{thm:theta3_weight1}
Let $\vartheta_3(\tau)^2 = \sum_{n=0}^{\infty} r_2(n)\, q^n$.
Then $\vartheta_3(\tau)^2 \in M_1(\Gamma_0(4), \chi_{-4})$,
the space of modular forms of weight $1$ for $\Gamma_0(4)$
with the non-principal character $\chi_{-4}(d) = \left(\frac{-4}{d}\right)$.
\end{theorem}

\begin{proof}
The theta series $\vartheta_3(\tau)^2 = \sum_{n\geq 0} r_2(n) q^n$
is holomorphic on $\mathfrak{H}$ and at all cusps (the coefficients $r_2(n)$
grow polynomially).  Under $T: \tau\to\tau+1$ it acquires $e^{2\pi i \cdot 0} = 1$
by Theorem~\ref{thm:T_transform} applied twice.
Under $S: \tau \to -1/\tau$ restricted to $\Gamma_0(4)$ (more precisely, under
all $\begin{pmatrix}a&b\\4c&d\end{pmatrix}\in\Gamma_0(4)$),
Theorem~\ref{thm:S_transform} gives the factor $\sqrt{-i\tau}^{\,2} = -i\tau$
(which is $(c\tau+d)^1$ for $c=0,d=1$ and corresponds to weight $1$).
The character $\chi_{-4}$ accounts for the behaviour at the cusp $\infty$.
A detailed verification using the explicit matrix entries can be found in \cite{Iwaniec1997}.
\end{proof}

\begin{corollary}[Dimension of $M_1(\Gamma_0(4),\chi_{-4})$]
The space $M_1(\Gamma_0(4),\chi_{-4})$ is one-dimensional, so
$\vartheta_3(\tau)^2$ is (up to a scalar) its unique element.
This uniqueness forces the explicit formula for $r_2(n)$ derived in Section~\ref{sec:quadratic}.
\end{corollary}

\subsection{The half-integer weight case: $\vartheta_3(\tau)$}

The function $\vartheta_3(\tau)$ itself is a modular form of weight $1/2$
for $\Gamma_0(4)$, in the sense of Shimura \cite{Shimura1973}.
The theta series $\vartheta_3(\tau)^k$ has weight $k/2$.
For even $k$ this is an integer-weight modular form for $\Gamma_0(4)$,
and its coefficients $r_k(n)$ can be expressed via Eisenstein series and cusp forms.

%% ============================================================
\section{Representations by Quadratic Forms}\label{sec:quadratic}
%% ============================================================

\subsection{Counting representations}

\begin{definition}
For $n \geq 0$ and $k \geq 1$, the \emph{representation number} $r_k(n)$ is
\[
  r_k(n) = \#\{(x_1,\ldots,x_k)\in\mathbb{Z}^k : x_1^2+\cdots+x_k^2 = n\}.
\]
Signed tuples and different orderings count separately.
\end{definition}

\begin{theorem}[Generating function]\label{thm:gen_rk}
$\displaystyle\vartheta_3(\tau)^k = \sum_{n=0}^{\infty} r_k(n)\, q^n$,
\quad $q = e^{\pi i\tau}$.
\end{theorem}

\begin{proof}
Expanding $(\sum_{m\in\mathbb{Z}} q^{m^2})^k$ by multinomial expansion,
the coefficient of $q^n$ is the number of tuples $(x_1,\ldots,x_k)$ with
$x_1^2+\cdots+x_k^2 = n$.
\end{proof}

\subsection{Jacobi's two-square formula}

\begin{theorem}[Two squares, Jacobi 1829]\label{thm:r2}
For $n \geq 1$,
\[
  r_2(n) = 4\sum_{d \mid n} \chi_{-4}(d)
  = 4\!\!\sum_{\substack{d \mid n \\ d \text{ odd}}} (-1)^{(d-1)/2}.
\]
Equivalently, $r_2(n) = 4(d_1(n) - d_3(n))$ where $d_j(n)$ is the number of
divisors of $n$ congruent to $j$ modulo $4$.
\end{theorem}

\begin{proof}
Since $M_1(\Gamma_0(4),\chi_{-4})$ is one-dimensional (Corollary after
Theorem~\ref{thm:theta3_weight1}), $\vartheta_3(\tau)^2$ is proportional to the
Eisenstein series $E_1(\tau,\chi_{-4}) = \sum_{n\geq1}(\sum_{d|n}\chi_{-4}(d))q^n$.
Comparing the coefficient of $q^1$: $r_2(1) = 4 = 4\cdot\chi_{-4}(1)$
fixes the proportionality constant to $4$. \qed
\end{proof}

\subsection{Jacobi's four-square formula}

\begin{theorem}[Four squares, Jacobi 1829]\label{thm:r4}
For all $n \geq 1$,
\begin{equation}\label{eq:r4}
  r_4(n) = 8\sum_{\substack{d \mid n \\ 4 \nmid d}} d.
\end{equation}
\end{theorem}

\begin{proof}
\textbf{Step 1 (Modular form identity).}
The function $f(\tau) = \vartheta_3(\tau)^4$ is a modular form of weight $2$
for $\Gamma_0(4)$.  By the dimension formula, $M_2(\Gamma_0(4))$ is spanned by
two Eisenstein series $E_2^{(0)}$ (at cusp $0$) and $E_2^{(\infty)}$ (at cusp $\infty$).
One computes:
\[
  \vartheta_3(\tau)^4
  = E_2^{(\infty)}(\tau) - 4\, E_2^{(0)}(\tau),
\]
where explicitly
$E_2^{(\infty)}(\tau) = 1 + 8\sum_{n\geq1}\bigl(\sum_{4\nmid d \mid n} d\bigr) q^n$.

\textbf{Step 2 (Initial conditions).}
The coefficient of $q^0$ is $1 = r_4(0)$ (the empty sum).
The coefficient of $q^1$: $r_4(1) = 8\cdot 1 = 8$
(the $8$ tuples $(\pm1,0,0,0)$ and permutations). \checkmark

\textbf{Step 3 (Uniqueness).}
Since $\vartheta_3(\tau)^4$ and the Eisenstein series have the same first two
Fourier coefficients and both lie in the same two-dimensional space,
they agree; extracting the $n$-th coefficient gives \eqref{eq:r4}.
A complete proof using Hurwitz's formula is in \cite{Grosswald1985}.
\end{proof}

\begin{corollary}[Lagrange's four-square theorem]
Every positive integer is a sum of four perfect squares.
\end{corollary}

\begin{proof}
For $n \geq 1$, the divisor sum in \eqref{eq:r4} includes $d = n$ (if $4\nmid n$)
or $d = n/4^k$ for the largest $k$ with $4^k \mid n$ (if $4\mid n$).
In either case the sum is positive, so $r_4(n) \geq 8 > 0$.
\end{proof}

%% ============================================================
\section{Poisson Summation and the Riemann Zeta Function}\label{sec:poisson}
%% ============================================================

\begin{theorem}[Poisson Summation Formula]\label{thm:poisson}
Let $f \in \mathcal{S}(\mathbb{R})$ (Schwartz class). Then
\[
  \sum_{n=-\infty}^{\infty} f(n) = \sum_{n=-\infty}^{\infty} \hat{f}(n),
  \qquad \hat{f}(\xi) = \int_{-\infty}^{\infty} f(x)\, e^{-2\pi i \xi x}\,dx.
\]
\end{theorem}

\begin{proof}
Let $F(x) = \sum_{n\in\mathbb{Z}} f(x+n)$; this is periodic with period $1$ and
lies in $C^\infty(\mathbb{R}/\mathbb{Z})$ by the rapid decay of $f$.
Its Fourier series is $F(x) = \sum_{m\in\mathbb{Z}} \hat{F}(m) e^{2\pi imx}$ with
$\hat{F}(m) = \int_0^1 F(x) e^{-2\pi imx}dx = \hat{f}(m)$.
Setting $x = 0$ gives $F(0) = \sum_n f(n) = \sum_m \hat{f}(m)$.
\end{proof}

\subsection{The theta function and $\zeta(s)$}

\begin{definition}
For $t > 0$, set
\[
  \theta(t) = \vartheta_3(it) = \sum_{n=-\infty}^{\infty} e^{-\pi n^2 t}.
\]
\end{definition}

\begin{theorem}[Functional equation of $\theta$]\label{thm:theta_fe}
\[
  \theta(t) = t^{-1/2}\,\theta(1/t), \qquad t > 0.
\]
\end{theorem}

\begin{proof}
Apply Poisson summation (Theorem~\ref{thm:poisson}) to $f_t(x) = e^{-\pi x^2 t}$.
Its Fourier transform is $\hat{f}_t(\xi) = t^{-1/2} e^{-\pi\xi^2/t}$.
Hence $\sum_n e^{-\pi n^2 t} = t^{-1/2} \sum_n e^{-\pi n^2/t}$.
\end{proof}

\begin{theorem}[Functional equation of $\zeta(s)$ via $\theta$]\label{thm:zeta_fe}
Let $\xi(s) = \tfrac{1}{2}s(s-1)\pi^{-s/2}\Gamma(s/2)\zeta(s)$.  Then
$\xi(s) = \xi(1-s)$.
\end{theorem}

\begin{proof}
\textbf{Step 1.}  From the Gamma integral,
$\pi^{-s/2}\Gamma(s/2) n^{-s} = \int_0^\infty t^{s/2-1} e^{-\pi n^2 t}\,dt$.
Summing over $n \geq 1$:
\[
  \pi^{-s/2}\Gamma(s/2)\zeta(s)
  = \int_0^\infty t^{s/2-1} \frac{\theta(t)-1}{2}\,dt.
\]

\textbf{Step 2.} Split the integral at $t=1$:
\[
  \int_0^1 + \int_1^\infty.
\]
In the integral over $(0,1)$, substitute $t \to 1/t$ and use
$\theta(1/t) = \sqrt{t}\,\theta(t)$ (Theorem~\ref{thm:theta_fe}):
\[
  \int_0^1 t^{s/2-1}\frac{\theta(t)-1}{2}\,dt
  = \int_1^\infty t^{(1-s)/2 - 1}\frac{\theta(t)-1}{2}\,dt
    + \frac{1}{s-1} - \frac{1}{s}.
\]

\textbf{Step 3.}  Combining both parts yields
\[
  \pi^{-s/2}\Gamma(s/2)\zeta(s)
  = \int_1^\infty \left(t^{s/2-1}+t^{(1-s)/2-1}\right)\frac{\theta(t)-1}{2}\,dt
    + \frac{1}{s(s-1)},
\]
which is symmetric under $s \leftrightarrow 1-s$, giving $\xi(s) = \xi(1-s)$.
\end{proof}

\begin{remark}
This is exactly the argument Riemann used in his 1859 paper
\textit{\"Uber die Anzahl der Primzahlen unter einer gegebenen Gr\"o\ss e} \cite{Riemann1859}.
The theta function functional equation is the key mechanism behind the symmetry.
\end{remark}

%% ============================================================
\section{Siegel Theta Functions}\label{sec:siegel}
%% ============================================================

\subsection{Lattices and theta series}

Let $\Lambda \subset \mathbb{R}^n$ be a lattice, i.e.\ a discrete subgroup of rank $n$.
The \emph{theta series} of $\Lambda$ is
\[
  \Theta_\Lambda(\tau) = \sum_{\mathbf{v} \in \Lambda} q^{\|\mathbf{v}\|^2/2},
  \qquad q = e^{2\pi i\tau}.
\]

\begin{example}
For $\Lambda = \mathbb{Z}^n$ with Euclidean norm,
$\Theta_{\mathbb{Z}^n}(\tau) = \vartheta_3(\tau)^n = \sum_{k=0}^\infty r_n(k)\, q^k$.
\end{example}

\subsection{Siegel's generalisation}

\begin{definition}[Siegel theta series]\label{def:siegel_theta}
For a positive-definite symmetric matrix $A \in M_n(\mathbb{Z})$ (the Gram matrix
of a quadratic form $Q(\mathbf{x}) = \mathbf{x}^T A \mathbf{x} / 2$) and
$\tau \in \mathfrak{H}$,
\[
  \Theta_A(\tau) = \sum_{\mathbf{x} \in \mathbb{Z}^n} e^{\pi i \tau\, Q(\mathbf{x})}.
\]
\end{definition}

\begin{theorem}[Modularity of $\Theta_A$]\label{thm:siegel_modular}
$\Theta_A(\tau)$ is a modular form of weight $n/2$ for a congruence subgroup
$\Gamma_0(N)$ depending on the level $N = \det(2A)$ of the lattice.
\end{theorem}

\begin{proof}[Proof sketch]
The transformation under $S: \tau \to -1/\tau$ follows from the multi-dimensional
Poisson summation formula applied to $f(\mathbf{x}) = e^{-\pi Q(\mathbf{x})/t}$ ($\tau = it$):
\[
  \Theta_A(-1/\tau) = \frac{\sqrt{-i\tau}^n}{\sqrt{\det A}}\,\Theta_{A^{-1}}(\tau).
\]
If $A$ is \emph{even} and \emph{unimodular} ($\det A = 1$, $A^{-1}=A^{-1}$ integral),
this simplifies to $\Theta_A(-1/\tau) = (-i\tau)^{n/2}\Theta_A(\tau)$,
showing $\Theta_A \in M_{n/2}(\mathrm{SL}_2(\mathbb{Z}))$.
For general $A$, one restricts to $\Gamma_0(N)$. See \cite{Serre1973} for details.
\end{proof}

\begin{example}[$E_8$ lattice]
The $E_8$ lattice is an even unimodular lattice in $\mathbb{R}^8$.
Its theta series $\Theta_{E_8}(\tau) \in M_4(\mathrm{SL}_2(\mathbb{Z}))$
equals the Eisenstein series $E_4(\tau) = 1 + 240\sum_{n\geq1}\sigma_3(n)q^n$.
Hence the number of vectors of norm $2k$ in $E_8$ equals $240\,\sigma_3(k)$.
\end{example}

%% ============================================================
\section{Connections to Elliptic Curves}\label{sec:elliptic}
%% ============================================================

\subsection{Uniformisation of elliptic curves}

Every elliptic curve $E$ over $\mathbb{C}$ is isomorphic (as a Riemann surface) to
a complex torus $\mathbb{C}/\Lambda_\tau = \mathbb{C}/(\mathbb{Z} + \tau\mathbb{Z})$
for some $\tau \in \mathfrak{H}$.  The uniformisation map
$\varphi: \mathbb{C}/\Lambda_\tau \to E(\mathbb{C})$
is given by the Weierstra\ss{} $\wp$-function.

\subsection{Expressing $\wp$ via theta functions}

\begin{theorem}[$\wp$ in terms of $\vartheta_1$]\label{thm:wp_theta}
\[
  \wp(z|\tau)
  = -\partial_z^2 \log\vartheta_1(z|\tau)
    + c(\tau),
\]
where $c(\tau)$ is a correction term independent of $z$ ensuring
$\wp$ has double poles at $z \in \Lambda_\tau$.
Explicitly,
\[
  \wp(z|\tau) = \left(\frac{\pi\vartheta_1'(0|\tau)}{\vartheta_1(z|\tau)}\right)^2
  - \frac{\pi^2\vartheta_1''(0|\tau)}{3\vartheta_1'(0|\tau)},
\]
where primes denote partial derivatives with respect to $z$.
\end{theorem}

\begin{proof}
Both sides have the same poles (double poles at $\Lambda_\tau$), the same
quasi-periodicity (the characteristic of an elliptic function), and equal constant
terms in the Laurent expansion at $z=0$.  By Liouville's theorem, any elliptic
function with specified poles and singular parts is determined up to a constant;
matching the constant fixes the formula. See \cite{WhittakerWatson1927}, Ch.~21.
\end{proof}

\subsection{The $j$-invariant}

\begin{theorem}[$j$-invariant via $\vartheta_2, \vartheta_3, \vartheta_4$]
\[
  j(\tau)
  = 256\,\frac{(\vartheta_2(\tau)^8 + \vartheta_3(\tau)^8 + \vartheta_4(\tau)^8)^3}
              {(\vartheta_2(\tau)\,\vartheta_3(\tau)\,\vartheta_4(\tau))^8}.
\]
\end{theorem}

The $j$-invariant classifies elliptic curves over $\mathbb{C}$ up to isomorphism,
and its expression via theta functions reflects the deep interplay between
the function theory on $\mathfrak{H}$ and the arithmetic of elliptic curves.

%% ============================================================
\section{Ramanujan-Style Identities and Mock Theta Functions}\label{sec:ramanujan}
%% ============================================================

\subsection{Rogers--Ramanujan identities}

\begin{theorem}[Rogers--Ramanujan, 1894/1913]\label{thm:RR}
For $|q| < 1$,
\begin{align}
  \sum_{n=0}^{\infty} \frac{q^{n^2}}{(q;q)_n}
  &= \prod_{n\geq1}\frac{1}{(1-q^{5n-1})(1-q^{5n-4})},
  \label{eq:RR1}\\
  \sum_{n=0}^{\infty} \frac{q^{n(n+1)}}{(q;q)_n}
  &= \prod_{n\geq1}\frac{1}{(1-q^{5n-2})(1-q^{5n-3})},
  \label{eq:RR2}
\end{align}
where $(q;q)_n = \prod_{k=1}^n(1-q^k)$ is the $q$-Pochhammer symbol.
\end{theorem}

\begin{proof}[Proof via Bailey chains]
Rogers (1894) discovered these identities via $q$-hypergeometric series manipulations.
The key technique is the \emph{Bailey chain}: starting from a Bailey pair $(a_n, b_n)$
satisfying $a_n = \sum_{r=0}^n \frac{b_r}{(q;q)_{n-r}(q;q)_{n+r}}$,
one generates new identities by the Bailey lemma.
Choosing $a_n = q^{n^2}$ produces \eqref{eq:RR1}.
Ramanujan rediscovered the identities independently around 1913.
MacMahon and others provided combinatorial interpretations:
\eqref{eq:RR1} states that the number of partitions of $n$ into parts
$\equiv \pm 1 \pmod{5}$ equals the number of partitions into parts
differing by at least $2$. A complete proof is in \cite{Andrews1998}.
\end{proof}

\subsection{Mock theta functions}

In his last letter to Hardy (1920), Ramanujan introduced \emph{mock theta functions},
mysterious $q$-series that ``mimic'' theta functions near the unit circle but
do not satisfy modular transformation laws in the usual sense.
The simplest third-order mock theta function is
\[
  f(q) = \sum_{n=0}^{\infty} \frac{q^{n^2}}{(-q;q)_n^2}.
\]

\begin{remark}[Status: unresolved for 80 years]
For eighty years, it remained unclear whether mock theta functions
could be placed within any systematic theory of modular forms.
Watson (1937) and others computed their asymptotics, but the
``modularity defect'' resisted explanation.
\end{remark}

\begin{theorem}[Zwegers 2002 -- Mock theta functions as harmonic Maass forms]\label{thm:zwegers}
\textup{(Proved by Zwegers \cite{Zwegers2002}.)}
Every Ramanujan mock theta function $h(\tau)$ can be completed to a
\emph{harmonic Maass form} $\widehat{h}(\tau) = h(\tau) + h^*(\tau)$,
where $h^*(\tau)$ is a non-holomorphic correction term (the ``shadow''),
such that $\widehat{h}$ transforms as a modular form of weight $1/2$
for a congruence subgroup of $\mathrm{SL}_2(\mathbb{Z})$.
\end{theorem}

\begin{proof}[Summary of Zwegers' construction]
For the third-order mock theta function $f(q)$ above,
Zwegers defines
\[
  h^*(\tau) = \sqrt{-i}\int_{-\bar\tau}^{i\infty}
  \frac{\overline{g(-\bar\tau')}}{\sqrt{\tau'+\tau}}\,d\tau',
\]
where $g(\tau)$ is the Mordell integral shadow of $f$.
One verifies that $\widehat{h} = h + h^*$ satisfies
$\widehat{h}\!\left(\frac{a\tau+b}{c\tau+d}\right) = (c\tau+d)^{1/2}\chi(a,b,c,d)\,\widehat{h}(\tau)$
for all $\begin{pmatrix}a&b\\c&d\end{pmatrix}\in\Gamma_0(576)$.
The holomorphic part $h(\tau)$ does not transform cleanly by itself precisely because
$h^*$ picks up the defect term.
This framework -- \emph{harmonic Maass forms} -- was subsequently developed into
a rich theory by Bruinier--Funke, Ono, Zagier and others.
\end{proof}

\begin{remark}
Zwegers' result resolved the mystery Ramanujan himself sensed when he wrote:
``\textit{I discovered very interesting functions recently which I call `Mock' $\vartheta$-functions...}''.
The connection between mock theta functions and harmonic Maass forms is now a central
topic in modern number theory, with applications to partitions, quadratic forms,
and even mathematical physics (string theory, moonshine).
\end{remark}

%% ============================================================
\section*{References}
%% ============================================================

\begin{thebibliography}{10}

\bibitem{Jacobi1829}
C.~G.~J.~Jacobi,
\textit{Fundamenta Nova Theoriae Functionum Ellipticarum},
K\"onigsberg, 1829.

\bibitem{WhittakerWatson1927}
E.~T.~Whittaker and G.~N.~Watson,
\textit{A Course of Modern Analysis}, 4th ed.,
Cambridge University Press, 1927.

\bibitem{Riemann1859}
B.~Riemann,
\"{U}ber die Anzahl der Primzahlen unter einer gegebenen Gr\"o\ss e,
\textit{Monatsberichte der K\"oniglich Preu\ss ischen Akademie der Wissenschaften zu Berlin},
pp.\ 671--680, 1859.

\bibitem{Mumford1983}
D.~Mumford,
\textit{Tata Lectures on Theta I, II, III},
Birkh\"auser, 1983--1991.

\bibitem{Serre1973}
J.-P.~Serre,
\textit{A Course in Arithmetic},
Springer, 1973.

\bibitem{Apostol1990}
T.~M.~Apostol,
\textit{Modular Functions and Dirichlet Series in Number Theory}, 2nd ed.,
Springer, 1990.

\bibitem{Iwaniec1997}
H.~Iwaniec,
\textit{Topics in Classical Automorphic Forms},
American Mathematical Society, 1997.

\bibitem{Shimura1973}
G.~Shimura,
On modular forms of half integral weight,
\textit{Annals of Mathematics} \textbf{97} (1973), 440--481.

\bibitem{Andrews1998}
G.~E.~Andrews,
\textit{The Theory of Partitions},
Cambridge University Press, 1998.

\bibitem{Grosswald1985}
E.~Grosswald,
\textit{Representations of Integers as Sums of Squares},
Springer, 1985.

\bibitem{Zwegers2002}
S.~Zwegers,
\textit{Mock Theta Functions},
PhD Thesis, Universiteit Utrecht, 2002.

\end{thebibliography}

\end{document}
